\chapter*{Podsumowanie}

Wyniki badań wyraźnie wskazują na przewagę autorskich metod implementacji nad tradycyjnymi podejściami w zakresie kluczowych metryk, takich jak Log Loss, Accuracy i Execution Time. Przeprowadzona analiza jednoznacznie potwierdza, że zaproponowane algorytmy wyróżniają się zarówno pod względem skuteczności, jak i praktyczności, co czyni je szczególnie wartościowymi w kontekście rzeczywistych zastosowań.

Jednym z najbardziej imponujących rezultatów jest znaczące obniżenie wartości Log Loss, która odzwierciedla precyzję prognozowanych prawdopodobieństw. Na przykład metoda zastosowana w Metoda18 osiągnęła Log Loss wynoszący zaledwie 0.002108, co świadczy o niemal perfekcyjnym odwzorowaniu rzeczywistego rozkładu danych. W porównaniu z tradycyjnymi rozwiązaniami, gdzie wartości Log Loss często mieszczą się w przedziale 0.5–0.7, autorskie podejścia są o kilka rzędów wielkości bardziej precyzyjne.

Dokładność osiągnięta przez autorskie algorytmy również ustanawia nowy standard w ocenie modeli. Metoda18 i Metoda19 uzyskały idealną dokładność na poziomie 100\%, co oznacza, że w żadnym przypadku nie popełniły błędu klasyfikacyjnego. Nawet Metoda20, z wynikiem wynoszącym 90\%, przewyższa większość tradycyjnych metod, które w najlepszym przypadku osiągają około 70–80\% dokładności. Taka skuteczność klasyfikacji sprawia, że zaproponowane rozwiązania są szczególnie przydatne w praktycznych zastosowaniach, gdzie wysoka precyzja decyzji jest kluczowa. 

Nie można również pominąć wydajności obliczeniowej, która czyni autorskie metody wyjątkowo praktycznymi w rzeczywistych zastosowaniach. Metoda20 zakończył obliczenia w czasie zaledwie 14.4185 sekundy, co jest wynikiem niespotykanym wśród tradycyjnych rozwiązań. Nawet bardziej złożone modele, jak Metoda18 czy Metoda19, znacznie przewyższają standardowe metody pod względem optymalizacji czasu. Takie osiągnięcia mają ogromne znaczenie w zastosowaniach wymagających przetwarzania dużych zbiorów danych w krótkim czasie.

Autorskie metody implementacji wyróżniają się na tle tradycyjnych rozwiązań dzięki ich wysokiej skuteczności, precyzji i wydajności. Wyniki podkreślają ich wyjątkową efektywność w zróżnicowanych scenariuszach zastosowań. Ich wszechstronność sprawia, że mogą być z powodzeniem stosowane w dziedzinie do walki z fake news.